\chapter{در جست‌وجوی نشانه‌ها؛ طرح موضوع و بررسی پیشینۀ پژوهش}
\label{chap:literature}
\section{ضرورت انجام پژوهش}
% Reconstructed and edited from physical pages 26--27 of the source PDF.
خودکشی یک مسئلۀ مهم سلامت عمومی با پیامدهای فردی و اجتماعی گسترده است. سازمان جهانی بهداشت در گزارش مورد استناد، شمار سالانۀ جان‌باختگان را در حدود ۷۰۰ هزار نفر توصیف می‌کند و بر پیشگیری و دسترسی به حمایت تأکید دارد \cite{ref1}. این پدیده نتیجۀ تعامل عوامل ژنتیکی، روان‌شناختی، اجتماعی و محیطی است؛ ازاین‌رو توضیح یا پیش‌بینی آن با یک نشانۀ منفرد کافی نیست \cite{ref3}. شناسایی زودهنگام می‌تواند فرصت مداخله ایجاد کند، اما نیازمند داده و ابزار ارزیابی معتبر است.

اهمیت مسئله در گروه‌های سنی جوان نیز در آمارهای سلامت عمومی دیده می‌شود \cite{ref2,ref5}. شکل‌گیری هویت، فشارهای تحصیلی و شغلی، شروع زندگی مستقل و محدودیت مهارت‌های مقابله‌ای می‌توانند بخشی از زمینۀ آسیب‌پذیری باشند. در کنار آن، ترس از قضاوت و انگ بیماری روانی ممکن است مراجعه به خدمات رسمی را دشوار کند. برخی افراد تجربه‌ها و احساسات خود را در محیط‌های آنلاین بیان می‌کنند؛ این رفتار امکان مطالعۀ نشانه‌های زبانی را فراهم می‌کند، اما مجوز تفسیر قطعی وضعیت بالینی هر نویسنده نیست \cite{ref6,ref7}.

شبکه‌های اجتماعی و یادداشت‌های متنی، داده‌هایی با حجم زیاد و ساختار ناهمگون ایجاد می‌کنند. پژوهش‌های زبان‌شناختی به الگوهایی مانند استفاده از ضمیر اول‌شخص، واژگان مرتبط با تنهایی، ناامیدی یا مرگ و تغییر لحن توجه کرده‌اند \cite{ref8,ref9}. روش‌های پردازش زبان طبیعی و یادگیری ماشین می‌توانند این الگوها را در چارچوب یک مجموعه‌دادۀ برچسب‌گذاری‌شده تحلیل کنند. مزیت آن‌ها مقیاس‌پذیری تحلیل است؛ محدودیت آن‌ها نیز وابستگی به تعریف برچسب، بافت، زبان و دادۀ آموزشی است \cite{ref10}.

مجموعه‌دادۀ استاندارد و نمایندۀ همۀ کاربران به‌سادگی در دسترس نیست. متن‌های شبکۀ اجتماعی با مصاحبۀ بالینی تفاوت دارند و نمونه‌گیری از یک بستر خاص می‌تواند سوگیری ایجاد کند. تفاوت فرهنگ، واژگان، سن، زمان و دسترسی کاربران به خدمات نیز بر تعمیم مدل اثر دارد. بنابراین ارزیابی باید همراه با تعریف روشنِ داده، کنترل تکرار، تفکیک کاربران یا متن‌های هم‌پوشان و بیان حدود کاربرد انجام شود. یک عدد دقت بالا، این مسائل را خودکار حل نمی‌کند.

مانع دیگر، حساسیت شدید اطلاعات سلامت روان است \cite{ref4}. افشای متن می‌تواند با انگ اجتماعی، تبعیض یا پیامدهای شخصی همراه باشد. قواعدی مانند \en{GDPR} و \en{HIPAA} نیز اهمیت تعیین هدف، دسترسی مجاز، نگهداری مسئولانه و رضایت را نشان می‌دهند. نیاز پژوهشگر به دادۀ متنوع باید با مسئولیت حفاظت از داده و محدودیت استفادۀ ثانویه سازگار شود. این نیاز، انگیزه‌ای برای بررسی روش‌هایی است که ضرورت جابه‌جایی دادۀ خام را کاهش می‌دهند.

یادگیری مشارکتی یکی از این مسیرهاست. در شکل‌های رایج آن، آموزش روی دستگاه یا سازمان محلی انجام و اطلاعاتی مانند به‌روزرسانی مدل یا خروجی پیش‌بینی به یک هماهنگ‌کننده فرستاده می‌شود. نوع اطلاعات اشتراکی به طراحی الگوریتم وابسته است. در روش این پایان‌نامه، اشتراک در سطح احتمال پیش‌بینی صورت می‌گیرد و دادۀ آموزش کلاینت‌ها برای آموزش یک مدل واحد به سرور منتقل نمی‌شود. این طراحی با میانگین‌گیری پارامترهای مدل یکسان نیست.

باقی ماندن دادۀ آموزش در محل، تنها یکی از لایه‌های حفاظت است. احتمال خروجی و به‌روزرسانی نیز می‌توانند اطلاعاتی حمل کنند؛ مهاجم فعال می‌تواند هویت جعل کند، وزن را تغییر دهد یا یک پیام معتبر را دوباره ارسال کند. در نتیجه، محرمانگی، اصالت، تمامیت و کنترل بازپخش باید به‌صورت جداگانه بررسی شوند. رمزگذاری همریخت، ابزار انجام محاسبه در دامنۀ رمز و امضای دیجیتال، ابزار بررسی اصالت محموله است؛ هیچ‌کدام جای تمام کنترل‌های سامانه را نمی‌گیرد.

ضرورت پژوهش حاضر از کنار هم قرار گرفتن این نیازها ناشی می‌شود: ارزیابیِ بدون نشت برای تحلیل متن، تعریف دقیق توزیع دادۀ کلاینت‌ها و سنجش یک روش احراز اصالت سازگار با تجمیع جمعی. هدف، گزارش یک چارچوب پژوهشیِ روشن و قابل ممیزی است. ادعای کاربرد بالینی، محرمانگی مطلق در برابر دارندۀ کلید یا اجرای اثبات دانش صفر، از آزمایش‌های انجام‌شده نتیجه گرفته نمی‌شود.

\section{نحوۀ انجام پژوهش}
پژوهش حاضر مسئلۀ تحلیل متن‌های مرتبط با خودکشی را در دو سطح کیفیت پیش‌بینی و حفاظت از تبادل مقادیر عددی بررسی می‌کند. در سطح نخست، یک خط‌لولۀ قابل‌بازتولید برای پاکسازی متن، یادگیری بازنمایی و مقایسۀ طبقه‌بندها تعریف می‌شود. در سطح دوم، چند مدل محلی روی زیرمجموعه‌های جداگانۀ داده آموزش می‌بینند و احتمال‌های خروجی آن‌ها با وزن تعداد نمونه‌های محلی تجمیع می‌شوند. توزیع طبقه‌بندی‌شدۀ داده، امکان بررسی اثر تعداد کلاینت‌ها را با حفظ تقریبی نسبت کلاس‌ها فراهم می‌کند.

\begin{figure}[htbp]
\centering
\begin{tikzpicture}[
  node distance=5mm,
  stage/.style={draw=blue!50!black,fill=blue!3,rounded corners=3pt,
    inner sep=6pt,font=\small},
  flow/.style={-{Latex[length=2.2mm]},thick,draw=blue!50!black}]
\node[stage] (partition) {\diagramtext{تخصیص بدون هم‌پوشانیِ آموزش به $K$ کلاینت\\
  حفظ تقریبی نسبت دو کلاس؛ هر ردیف دقیقاً در یک بخش محلی}};
\node[stage,below=of partition] (local) {\diagramtext{کلاینت $k$: یادگیری \en{Word2Vec} و \en{XGBoost} بر $D_k$\\
  واژگان و مدل محلی از اعتبارسنجی و آزمون یاد نمی‌گیرند}};
\node[stage,below=of local] (infer) {\diagramtext{برای متن ارزیابی $x$، هر کلاینت احتمال $p_k(x)$ را برمی‌گرداند\\
  پارامترهای درخت یا مختصات بازنمایی با هم میانگین گرفته نمی‌شوند}};
\node[stage,below=of infer] (aggregate) {\diagramtext{تجمیع احتمال‌ها با وزن اندازۀ دادۀ محلی\\
  $p_{\mathrm{ens}}(x)=\sum_{k=1}^{K}\frac{|D_k|}{\sum_j|D_j|}p_k(x)$}};
\node[stage,below=of aggregate] (evaluate) {\diagramtext{انتخاب آستانه بر اعتبارسنجی؛ ارزیابی ثابت بر آزمون مشترک\\
  خروجی‌های یادگیری در نوت‌بوک به‌صورت آشکار تجمیع شده‌اند}};
\foreach \a/\b in {partition/local,local/infer,infer/aggregate,aggregate/evaluate}
  \draw[flow] (\a) -- (\b);
\end{tikzpicture}

\caption{جریان مشارکتیِ این آزمایش: آموزش محلی و میانگین وزن‌دار احتمال‌ها با وزن تعداد نمونه}
\label{fig:federated-architecture}
\end{figure}

برای مقایسۀ متمرکز، چهار خانوادۀ جنگل تصادفی، درخت تصمیم، \en{XGBoost} و رأی‌گیری بیز ساده بررسی می‌شوند. مدل و آستانه بر اساس دادۀ اعتبارسنجی انتخاب می‌شوند و دادۀ آزمون تا تثبیت انتخاب نهایی کنار گذاشته می‌شود. یادگیری \en{Word2Vec} درون خط‌لوله قرار دارد؛ در اعتبارسنجی متقابل، هر بار فقط بخش آموزشیِ همان تقسیم برای یادگیری واژگان و بردارها استفاده می‌شود. این طراحی از ورود واژگان و الگوهای مجموعۀ ارزیابی به یادگیری بازنمایی جلوگیری می‌کند.

برای آزمایش مشارکتی، سناریوهای ۱۰، ۲۰، ۵۰ و ۱۰۰ کلاینت روی مجموعۀ آموزش ثابت تعریف شده‌اند. هر کلاینت، \en{Word2Vec} و طبقه‌بند خود را فقط روی دادۀ محلی آموزش می‌دهد. پرسش‌های اعتبارسنجی و آزمون برای همه یکسان‌اند و سرور احتمال کلاس مثبت را با وزن اندازۀ دادۀ محلی ترکیب می‌کند. این پیاده‌سازی با میانگین‌گیری وزن درخت‌ها یا الگوریتم \en{FedAvg} یکسان نیست؛ واحد اشتراک در آن، احتمال پیش‌بینی است.

لایۀ امنیتی نیز به‌صورت جداگانه بررسی می‌شود. رمزگذاری متقارن \en{AES-GCM} و تسهیم راز شامیر برای توضیح حفاظت از ورودی و مدیریت کلید به کار می‌روند. سپس طرح‌های همریخت جمعی و ضربی مقایسه می‌شوند و ساختار \en{SBHE-Benaloh} برای احراز اصالت محموله و تجمیع مجموع وزن‌دار معرفی می‌شود. امضای دیجیتال در کنار متن رمز قرار می‌گیرد و جزئی از ضرب جبری متن‌های رمز نیست. آزمایش زمانی رمزنگاری، آزمایش طبقه‌بندی و اجرای یک سامانۀ سرتاسری سه مفهوم متفاوت‌اند؛ نتایج هر کدام با حدود اعتبار خود گزارش می‌شوند.

\section{پیشینۀ پژوهش}
\subsection{تشخیص خطر خودکشی با تحلیل متن و یادگیری ماشین}
روش‌های سنتی ارزیابی خطر خودکشی به مصاحبۀ بالینی، پرسش‌نامه، مشاهده و قضاوت متخصص متکی‌اند. ارزش این روش‌ها در دسترسی به زمینۀ فردی و تفسیر تخصصی است، اما زمان ارزیابی، دسترسی محدود به متخصص و نبود مشاهدۀ پیوستۀ رفتار، کاربرد آن‌ها را در مقیاس وسیع دشوار می‌کند. تحلیل محاسباتی متن برای تکمیل مطالعۀ نشانه‌ها مطرح شده است؛ جایگاه آن باید از تشخیص بالینی مستقل تفکیک شود.

جی و همکاران در یک مرور جامع، روش‌های یادگیری ماشین و یادگیری عمیق برای تشخیص افکار خودکشی را بررسی کردند. منابع دادۀ مرورشده شامل پرسش‌نامه‌ها، پرونده‌های سلامت، یادداشت‌ها و محتوای تولیدشده در محیط آنلاین بودند. این تنوع نشان می‌دهد که «تشخیص خودکشی از متن» یک وظیفۀ یکنواخت با برچسب و شرایط ثابت نیست؛ منبع داده و تعریف هدف باید در کنار معیار مدل گزارش شوند \cite{ref11}.

هکلر و همکاران نیز در مرور نظام‌مند خود به تمرکز پژوهش‌ها بر دادۀ شبکه‌های اجتماعی، افزایش استفاده از یادگیری عمیق و حرکت از ویژگی‌سازی دستی به استخراج خودکار ویژگی توجه کردند. در عین حال، وابستگی به دادۀ خوداظهاری، محدودیت تعمیم و نیاز به شناسایی نشانه‌هایی فراتر از گزارش مستقیم فرد، از چالش‌های مطرح‌شده بودند \cite{ref12}. بنابراین رشد توان بازنمایی مدل، مسئلۀ اعتبار برچسب و نمایندگی نمونه را حذف نمی‌کند.

انجمن‌هایی مانند \en{Reddit} به دلیل وجود متن‌های طولانی و خوداظهاری، منبع متداول پژوهش شده‌اند. مطالعات این حوزه معمولاً زنجیره‌ای از گردآوری، تعریف برچسب، پاکسازی، استخراج ویژگی، آموزش و ارزیابی را اجرا می‌کنند \cite{ref13}. مرورهای مرتبط بر نقش الگوهای زبانی و امکان مطالعۀ نشانه‌های عاطفی تأکید دارند \cite{ref14}. با این حال، ارزیابی بر یک بستر خاص، گواه موفقیت در زبان یا جامعۀ دیگری نیست.

بخش مهمی از این ادبیات، داده را در محیط متمرکز در اختیار پژوهشگر قرار می‌دهد. این فرض برای توسعۀ اولیه ساده است، اما در کاربرد سلامت روان، محدودیت مالکیت، رضایت و محرمانگی ایجاد می‌کند. متون ممکن است اطلاعات شخصی، پزشکی و عاطفی را هم‌زمان حمل کنند. ازاین‌رو افزایش دقت، بدون توجه به محل پردازش و نوع اطلاعات مبادله‌شده، مسئلۀ حفاظت از داده را کامل پاسخ نمی‌دهد.

\subsection{یادگیری مشارکتی در حوزۀ سلامت}
مک‌ماهان و همکاران، آموزش مدل از داده‌های توزیع‌شده با به‌روزرسانی‌های محلی و تجمیع را به‌صورت نظام‌مند بررسی کردند \cite{ref15}. در این چارچوب، دادۀ خام لازم نیست برای ساخت یک مخزن آموزشی واحد منتقل شود. روش‌های مشخص، در نوع مدل، تعداد دور، شیوۀ تجمیع و فرض توزیع داده متفاوت‌اند. بنابراین هر پیاده‌سازی که چند کلاینت دارد الزاماً همان الگوریتم \en{FedAvg} نیست.

در سلامت دیجیتال، داده‌ها معمولاً در بیمارستان‌ها، مراکز درمانی و سازمان‌های مستقل نگهداری می‌شوند. ریک و همکاران، جدایی این منابع و موانع حقوقی، اخلاقی و سازمانی انتقال داده را از انگیزه‌های مهم یادگیری مشارکتی دانسته‌اند \cite{ref16}. امکان همکاری بدون متمرکزسازیِ متن خام، یک مزیت بالقوه است؛ در مقابل، سازگاری قالب، کیفیت متفاوت داده و محدودیت ارتباط باید در طراحی لحاظ شوند.

کایروز و همکاران در بررسی جامع خود به ارتباطات، ناهمگونی، حریم خصوصی، امنیت و مقیاس‌پذیری پرداخته‌اند \cite{ref17}. محلی ماندن آموزش می‌تواند اصل حداقل‌سازی دادۀ انتقالی را تقویت کند، اما به‌روزرسانی یا خروجی محلی نیز ممکن است اطلاعات حساس حمل کند. حملۀ استنتاج عضویت، وارون‌سازی مدل یا تحلیل خروجی، نمونه‌هایی از علت نیاز به مدل تهدیدِ صریح‌اند.

کاربرد یادگیری مشارکتی در سلامت روان، از جمله تشخیص افسردگی با مدل‌های زبانی و پست‌های شبکۀ اجتماعی، نیز بررسی شده است. مطالعۀ مورد ارجاع، داده‌های مستقل و غیرهم‌توزیع را در چارچوبی برای همکاری منابع داده مطرح می‌کند \cite{ref18}. این پیشینه، انگیزه‌ای برای سنجش روش‌های توزیع‌شده فراهم می‌کند؛ نتیجۀ یک وظیفۀ سلامت روان را نمی‌توان بی‌واسطه به وظیفۀ دیگر یا به جمعیت بالینی تعمیم داد.

برای مسئلۀ حاضر، نوع تقسیم کلاینت اهمیت ویژه دارد. تعداد کلاینت، اندازۀ دادۀ محلی و نسبت کلاس‌ها باید هم‌زمان گزارش شوند. تقسیم تصادفیِ طبقه‌بندی‌شده برای مطالعۀ اثر اندازه مناسب است، اما تفاوت واقعی کاربران، موضوع، زبان و زمان را شبیه‌سازی نمی‌کند. بخش مشارکتی این پایان‌نامه، همین محدودۀ آزمایش را روشن نگه می‌دارد.

\subsection{رمزگذاری همریخت و محاسبات حفظ‌کنندۀ حریم خصوصی}
ویژگی اصلی رمزگذاری همریخت، سازگاری برخی عملیات بر متن رمز با عملیات مشخص بر پیام است. پاییه یک نمونۀ مهمِ همریختی جمعی است که امکان ساخت متن رمزِ مجموع را فراهم می‌کند \cite{ref19}. این تناسب برای آمارگیری و تجمیع خروجی‌های عددی مفید است، اما نوع عملیات و فضای پیام باید با تابع مورد نظر سازگار باشند.

معرفی رمزگذاری کاملاً همریخت، چشم‌انداز انجام توابع پیچیده‌تر در دامنۀ رمز را گسترش داد \cite{ref20}. با این حال، قابلیت نظریِ محاسبۀ عمومی با هزینۀ یک جمع وزن‌دار در یک طرح جزئی همریخت یکسان نیست. انتخاب ابزار باید به تابع، پارامتر امنیت، نمایش عددی و هزینۀ مورد قبول توجه کند؛ پیچیده‌تر بودن ساختار، به‌تنهایی مزیت عملی برای هر کاربرد ایجاد نمی‌کند.

مرورهای کاربرد رمزگذاری همریخت در یادگیری ماشین پزشکی و زیست‌اطلاعات، ظرفیت این روش‌ها را برای طبقه‌بندی، تحلیل دادۀ ژنتیکی و پردازش اطلاعات حساس بررسی کرده‌اند \cite{ref21}. مرورهای حوزۀ سلامت نیز تفاوت پردازش بر متن رمز با الگوی «ابتدا رمزگشایی، سپس محاسبه» را برجسته می‌کنند \cite{ref22}. در عین حال، مشخص بودن دارندۀ کلید خصوصی و اطلاعاتی که از خروجی نهایی آشکار می‌شود، برای تفسیر محرمانگی ضروری است.

\subsection{تجمیع امن و ترکیب آن با یادگیری مشارکتی}
در تجمیع امن، هدف آن است که تابعی از مشارکت‌ها محاسبه شود و دسترسی به مقادیر منفرد مطابق مدل اعتماد محدود بماند. پروتکل بوناویتز و همکاران نمونه‌ای مهم در یادگیری ماشین حفظ‌کنندۀ حریم خصوصی است \cite{bonawitz}. این ادبیات نشان می‌دهد که «سرور هماهنگ‌کننده است» و «سرور باید تمام خروجی‌های منفرد را بداند» دو فرض یکسان نیستند.

روش‌های تجمیع امن ممکن است از تسهیم راز، پوشان‌سازی، رمزگذاری یا ترکیب آن‌ها استفاده کنند. تحمل خروج عضو، تبانی، پایداری وضعیت و اطلاعات آشکارشده از مجموع، در هر طرح متفاوت است. ازاین‌رو یک امضای معتبر یا یک ضرب همریخت، جایگزین تعریف کامل پروتکل و مدل تهدید نمی‌شود. در پیاده‌سازی این پایان‌نامه، سرور کلید خصوصی بنالوح را نگه می‌دارد و این تفاوت با طرح‌های دارای حفاظت در برابر سرور، صریحاً بیان می‌شود.

\subsection{جمع‌بندی پیشینه و شکاف پژوهشی}
سه جریان اصلی در پیشینه دیده می‌شوند: تحلیل زبانیِ نشانه‌های خطر، همکاری منابع داده بدون متمرکزسازی آموزش و محاسبۀ رمزگذاری‌شده یا تجمیع امن. جریان نخست به نمایش و کیفیت پیش‌بینی، جریان دوم به محل یادگیری و ناهمگونی و جریان سوم به اطلاعات قابل مشاهده و مدل اعتماد می‌پردازد. ارزیابی یک سامانۀ ترکیبی باید این سه دامنه را به هم مرتبط کند، اما معیارهای آن‌ها را با هم اشتباه نگیرد.

\begin{figure}[htbp]
\centering\begin{tikzpicture}[
  node distance=5mm,
  stage/.style={draw=teal!60!black,fill=teal!3,rounded corners=3pt,
    inner sep=6pt,font=\small},
  flow/.style={-{Latex[length=2.2mm]},thick,draw=teal!60!black}]
\node[stage] (nlp) {\diagramtext{تحلیل متن سلامت روان\\
  تعریف برچسب، نمایش متن، جداسازی داده و اعتبار ارزیابی}};
\node[stage,below=of nlp] (fl) {\diagramtext{یادگیری مشارکتی در سلامت\\
  آموزش محلی، توزیع کلاینت‌ها، ناهمگونی و روش تجمیع}};
\node[stage,below=of fl] (crypto) {\diagramtext{محاسبۀ رمزگذاری‌شده و تجمیع امن\\
  عملیات مجاز، مدل اعتماد، دارندۀ کلید و اطلاعات آشکارشده}};
\node[stage,below=of crypto] (thesis) {\diagramtext{مسیر ارزیابی این پایان‌نامه\\
  مقایسۀ متمرکز، مجموعۀ مدل‌های محلی و بنچمارک مستقل رمزنگاری\\
  اتصال سرتاسری نیازمند سنجش نمایش عددی، ظرفیت و سیاست کلید است}};
\draw[flow] (nlp) -- (fl);
\draw[flow] (fl) -- (crypto);
\draw[flow] (crypto) -- (thesis);
\end{tikzpicture}

\caption{رابطۀ سه حوزۀ پیشینه با مسیرهای ارزیابی پژوهش حاضر}
\label{fig:literature-overview}
\end{figure}

نیاز پژوهشیِ مورد توجه، تعریف یک مسیر قابل ممیزی است: مقایسۀ مدل با جداسازی درست داده، تعریف دقیق تخصیص کلاینت، گزارش احتمال تجمیعی و بررسی احراز اصالتِ سازگار با همریختی. این صورت‌بندی، ادعای نخستین ترکیب این حوزه‌ها در تمام ادبیات نیست. سهم این مخزن در پیاده‌سازی، مستندسازی پیش‌شرط‌ها و گزارش نتایج مشخصِ آزمایش‌ها ارزیابی می‌شود.

هدف معماری، کاهش نیاز به انتقال دادۀ آموزش و بررسی حفاظت از مشارکت‌های عددی است. رسیدن به یک سرویس امنِ سرتاسری مستلزم سازگاری وزن، ظرفیت پیام، کوانتیزه‌سازی، مدیریت هویت و سیاست رمزگشایی است. نتیجۀ طبقه‌بندی، نتیجۀ بنچمارک رمزنگاری و تحلیل امنیتی، در ادامۀ پایان‌نامه با همین تفکیک گزارش می‌شوند.


\section{صورت‌بندی پرسش‌ها و دامنۀ اعتبار}
پرسش نخست این است که با یک روش ارزیابی بدون نشت داده، کدام طبقه‌بند و آستانه بر دادۀ اعتبارسنجی مناسب‌تر است و عملکرد آن انتخابِ تثبیت‌شده در آزمون چه اندازه است. پرسش دوم اثر کاهش اندازۀ دادۀ محلی بر کیفیت پیش‌بینی تجمیع‌شده است. پرسش سوم، هزینه و محدودیت‌های پارامترهای یک پروتکل احراز اصالت و تجمیع همریخت است. هیچ‌یک از این پرسش‌ها به‌تنهایی اثبات نمی‌کند که طبقه‌بند برای تصمیم‌گیری بالینی آماده است یا سرورِ مالک کلید خصوصی از نظر رمزنگاری قادر به رمزگشایی خروجی منفرد نیست.

مدل تهدید، نقش مالک کلید خصوصی، شرط ظرفیت فضای پیام و تفاوت میان صحت امضا و صحت محتوای پیش‌بینی باید صریح بمانند. همچنین، آزمایش هم‌توزیع، جایگزین بررسی ناهمگونی واقعی کاربران، زبان‌ها، سازمان‌ها یا بازه‌های زمانی نیست. این محدودیت‌ها مبنای پیشنهادهای توسعۀ آتی در پایان پژوهش قرار می‌گیرند.
